\section{Introduction and Justification}

The $L_0$ regularization and its generalizations for group sparsity are fundamental tools in machine learning and signal processing, enabling variable selection and the interpretability of complex mathematical models. Due to the nonconvex and combinatorial nature of the $L_0$ penalty, efficiently solving these large-scale problems remains an open challenge in optimization. Other regularizations, such as $L_1$ (Lasso) and $L_2$ (ridge) hybrids with $L_0$, are also of recent interest for their potential to avoid overfitting while maintaining the benefits of pure $L_0$ regularization \cite{fastselect}.

During the ongoing PhD research in Brazil, the central focus has been the development of highly efficient first-order heuristics for these problem classes. Specifically, we have developed algorithms based on the Nonmonotone Spectral Proximal Gradient (NSPG) combined with Coordinate Descent (CD) and combinatorial local search methods. These nonmonotone methods act efficiently as a preprocessing step (rough support identification) for state-of-the-art algorithms (e.g., PGCCD \cite{fastselect}), overcoming dimensionality challenges in high complexity and parameter correlation environments. The implementation was developed using the Julia language, which allows high performance and flexibility in developing optimization algorithms.

Despite the high computational efficiency of first-order methods (spectral), their capacity to act as global preprocessors and identify good regions of the search space could benefit substantially from second-order approaches. Enhancing this stage with structured proximal operators allows the algorithm to get much closer to the correct manifold, ensuring that the subsequent coordinate descent and combinatorial algorithms have an easier time performing local refinement and finding overall better solutions. It is in this context that the proposed international collaboration is founded.

\subsection{Justification for the Host Institution}
The research group led by Prof. Dominique Orban at Polytechnique Montréal (GERAD) possesses internationally recognized leadership and excellence in the development of scalable nonsmooth and second-order optimization methods. Recent works from the group, such as the Levenberg-Marquardt method for nonsmooth regularized least squares \cite{aravkin2024levenberg} and scalable adaptive cubic regularization methods \cite{dussault2024scalable}, represent the state-of-the-art in the robust resolution of complex optimization problems. Additionally, recent analyses on the complexity of trust-region methods for nonsmooth optimization \cite{leconte2023complexity} offer a solid theoretical foundation for extensions into discontinuous regularizations of interest. Furthermore, the group has demonstrated a strong interest in the study of nonmonotone algorithmic behavior and the development of alternative proximal operators \cite{Leconte_2024}, which closely aligns with the applicant's current research trajectory.

Beyond theoretical alignment, the host group is the main developer of \texttt{JuliaSmoothOptimizers} (JSO), a widely used ecosystem for optimization in Julia. Immersion in this development environment will foster a unique synergy, allowing not only the theoretical advancement of the research initiated in Brazil but also its conversion into impactful, globally used software. Thus, the research internship abroad will bring a substantial contribution that could not be fully obtained solely at the home institution.
